\section{The unconditional lower bound from all zero jets}

This section identifies the measure which every subsequential cutoff limit
must dominate.  The argument uses no location theorem for the zeros.
We state its analytic inputs explicitly so that the conclusion does not
inherit a Riemann-hypothesis assumption from an earlier Gram estimate.

Let $\chi$ be a fixed primitive Dirichlet character of conductor $q$.
Put $R=\Lambda/q$, $c=\Lambda R$, and $L=\log\Lambda$.
The raw Hilbert space is the parity-$\chi(-1)$ subspace of
$L^2(\mathbb R,du)$, normalized by its positive-half-line norm.
Write $P$ for time restriction to $[-\Lambda,\Lambda]$ and $B$ for
Fourier restriction to $[-R,R]$, with Fourier phase $e^{-2\pi iu\xi}$.
The selected raw space is spanned by the time-supported prolate
eigenfunctions of $PBP$ with the prescribed parity and full indices
$n\leq4c$.  Its projection is $\Pi_\Lambda$.
For the principal character intersect this space with
$\ker I\cap\ker\operatorname{ev}_0$, where
$I\psi=\int_0^\infty\psi(u)\,du$.

The theta map into $\mathcal H=L^2(\mathbb R,dt)$ is
\[
 E_\chi\psi(t)=e^{t/2}\sum_{n\geq1}\chi(n)\psi(ne^t)
 \quad(\chi\ne1),
\]
and in the principal sector it has the additional term
$-I\psi\,e^{-t/2}$.
Choose a raw orthonormal basis of the selected space, let $E_\Lambda$
be its theta column operator, and put
\[
 G_\Lambda=E_\Lambda^*E_\Lambda,\qquad
 Z_\Lambda=E_\Lambda G_\Lambda^{-1/2},\qquad
 Q_\Lambda=Z_\Lambda Z_\Lambda^*,\qquad
 S_L=1_{[-L,L]}.
\]
The vectors $z_i=Z_\Lambda e_i$ are orthonormal in $\mathcal H$.
We use $\widehat f(\tau)=\int f(t)e^{i\tau t}\,dt$, with Plancherel
measure $d\tau/(2\pi)$, and $C_a=\mathcal F^{-1}a\mathcal F$.
Set $D_L=S_L-S_LQ_\Lambda S_L$.

The required analytic inputs are the following.
First, the selected concentration eigenvalues have a common positive
lower bound $\nu$ for large $\Lambda$, by Lemma~\ref{lem:raw-estimates}.
Second, Proposition~\ref{prop:gram-control} gives, for every fixed $B_0>0$,
\begin{equation}\label{eq:alljet-gram-input}
 \|G_\Lambda^{-1}\|\leq
 C_{\chi,B_0}\Lambda^2(\log\Lambda)^{-2B_0}.
\end{equation}
Third, Lemma~\ref{lem:theta-identities} gives the lower support gap for band
vectors, and Mellin summation gives
\[
 \int_{\mathbb R}E_\chi\psi(t)e^{(s-1/2)t}\,dt
 =L(s,\chi)\int_0^\infty\psi(u)u^{s-1}\,du
 \quad(0<\Re s<1),
\]
with the principal mean subtraction and constraints just stated.
The displayed integrals and their derivatives converge for each selected
vector.  Finally, when replacing the positive contraction by the
uncompressed projection, one uses
\begin{equation}\label{eq:alljet-trace-input}
 \operatorname{Tr}\bigl((Q_\Lambda-S_LQ_\Lambda S_L)V(h)\bigr)=o(1)
 \qquad(h\in C_c^\infty(\mathbb R)).
\end{equation}
Here $V(h)$ is logarithmic convolution; Section~\ref{sec:trace-transfer}
proves this last estimate. These are separate inputs;
the lower bound below only needs the first three until its last sentence.

\begin{theorem}\label{thm:allzero-lower}
Assume the analytic inputs above, with no hypothesis on zero location.
For a nontrivial zero $\rho=\beta+i\gamma$ write
$\alpha=\beta-1/2$ and let $m_\rho$ be its multiplicity.
Define the positive measure
\[
 d\mathfrak b_\chi(\tau)=
 \sum_{\alpha=0}m_\rho\,\delta_\gamma(d\tau)
 +\sum_{\alpha\ne0}m_\rho
 \frac{|\alpha|\,d\tau}{\pi((\tau-\gamma)^2+\alpha^2)}.
\]
The second sum is locally finite as a measure and is integrable against
every Schwartz function.  For every $a\in\mathcal S(\mathbb R)$,
\begin{equation}\label{eq:allzero-lower}
 \liminf_{\Lambda\to\infty}\operatorname{Tr}(D_LC_aC_a^*)
 \geq\int_{\mathbb R}|a(\tau)|^2\,d\mathfrak b_\chi(\tau).
\end{equation}
In particular every locally finite vague limit of the positive defect
measures dominates $\mathfrak b_\chi$.
If \eqref{eq:alljet-trace-input} also holds, the same inequality holds
for the corresponding uncompressed renormalized trace on convolution
squares of compactly supported smooth logarithmic functions.
\end{theorem}

\begin{proof}
Fix a zero $s_0=\beta+i\gamma$, and an integer $0\leq j<m_{s_0}$.
The lower theta-tail functional is
\[
 \ell_{s_0,j,\Lambda}(\psi)=
 \int_{-\infty}^{-L}E_\chi\psi(t)
 e^{(s_0-1/2)t}(t+L)^j\,dt.
\]
In a nonprincipal sector its raw kernel is
\[
 H_{s_0,j,\Lambda}(u)=
 u^{s_0-1}\sum_{n>\Lambda u}\chi(n)n^{-s_0}
                 \bigl(\log(\Lambda u/n)\bigr)^j.
\]
For $u<\Lambda^{-1}$ the zero and its first $j$ derivatives make this
kernel vanish.  For $u\geq\Lambda^{-1}$, bounded partial sums of
$\chi$ and one Abel summation give
\begin{equation}\label{eq:offaxis-tail-kernel}
 |H_{s_0,j,\Lambda}(u)|\leq C_{s_0,j,\chi}\Lambda^{-\beta}u^{-1}.
\end{equation}
Indeed the integral of the absolute derivative of
$v^{-s_0}(\log(\Lambda u/v))^j$ for $v>\Lambda u$ is
$O((\Lambda u)^{-\beta})$, and the endpoint has the same size.
The kernel calculation is justified first with a positive lower
cutoff on $e^t$; Abel summation bounds the omitted contribution by
$O(\varepsilon^\beta(1+|\log\varepsilon|)^{j+1})$ for a bounded
compact raw vector.  This tends to zero because $\beta>0$.

For the principal character the regularized kernel is instead
\[
 \left.(\partial_s+L)^j
 \left[u^{s-1}\left(\zeta(s)-\sum_{n\leq\Lambda u}n^{-s}\right)
             +\frac{\Lambda^{1-s}}{1-s}\right]\right|_{s=s_0}.
\]
Euler summation with the bounded discrepancy $\lfloor v\rfloor-v$
proves \eqref{eq:offaxis-tail-kernel} on $u\geq\Lambda^{-1}$.
On the remaining interval it is $O(\Lambda^{1-\beta})$.
Thus in both cases the functional is bounded on raw $L^2$ and
\[
 \|\ell_{s_0,j,\Lambda}(1-P)\|^2
                  \leq C_{s_0,j,\chi}\Lambda^{-2\beta-1}.
\]
The band support identity annihilates this functional on the band
range.  In the principal case the band extension must have value
zero at the origin.  For a constrained selected vector $\psi$, the
extension $w=B(PBP)^{-1}\psi$ satisfies $Pw=\psi$, hence $w(0)=0$.
Its extension operator has norm at most $\nu^{-1}$.
Consequently summing over a raw orthonormal selected basis gives,
in either character type,
\[
 \sum_i|\ell_{s_0,j,\Lambda}(\psi_i)|^2
                   \leq C_{s_0,j,\chi}\Lambda^{-2\beta-1}.
\]
In the nonprincipal sector one can obtain the same assertion directly
from $\Pi_\Lambda\leq\nu^{-1}PBP$.
The constants absorb the fixed powers of $\nu^{-1}$.
After multiplying by the inverse Gram square root and using
\eqref{eq:alljet-gram-input}, this yields
\begin{equation}\label{eq:offaxis-tail-normalized}
 \sum_i|\ell_{s_0,j,\Lambda}(z_i)|^2
 \leq C_{s_0,j,\chi,B_0}\Lambda^{-2\alpha}L^{-2B_0}.
\end{equation}

The full differentiated Mellin integral at $s_0$ vanishes.
The upper target tail is zero, including in the constrained principal
sector.  Therefore each window integral is the negative of its lower
tail.  For $\alpha<0$ use the window vectors
\[
 v_{s_0,j,L}(t)=1_{[-L,L]}(t)
            e^{(\alpha-i\gamma)(t+L)}(t+L)^j.
\]
Multiplication by $\Lambda^{\alpha}$ in this normalization cancels
the power $\Lambda^{-2\alpha}$ in
\eqref{eq:offaxis-tail-normalized}; hence
$\|Q_\Lambda v_{s_0,j,L}\|^2=O(L^{-2B_0})$.
For $\alpha>0$ use
\[
 v_{s_0,j,L}(t)=1_{[-L,L]}(t)
            e^{(\alpha-i\gamma)(t-L)}(t-L)^j.
\]
Expanding $(t-L)^j=((t+L)-2L)^j$ gives
\[
 \|Q_\Lambda v_{s_0,j,L}\|^2
                  =O(\Lambda^{-4\alpha}L^{2j-2B_0})=o(1).
\]
At $\alpha=0$ use the normalized vectors
\[
 v_{\gamma,j,L}(t)=1_{[-L,L]}(t)e^{-i\gamma t}
       \frac{((t+L)/(2L))^j}{\sqrt{2L}}.
\]
Equation \eqref{eq:offaxis-tail-normalized} again shows that their
$Q_\Lambda$-projections tend to zero.

Keep an arbitrary finite packet of zeros, including all their jets.
After translation of the upper endpoint to zero, the positive-$\alpha$
vectors converge in $L^2$ to
$1_{t\leq0}e^{(\alpha-i\gamma)t}t^j$.
The negative-$\alpha$ vectors converge similarly at the lower endpoint
to $1_{t\geq0}e^{(\alpha-i\gamma)t}t^j$.
Each finite endpoint Gram matrix is positive definite: a finite
exponential polynomial vanishing on an interval is identically zero,
and distinct exponents with their derivative orders are linearly
independent.  Opposite endpoints are asymptotically orthogonal.
Each endpoint packet has inner products $O(L^{-1/2})$ with a
critical-line packet, because the latter has bounded normalized
polynomial amplitude while an endpoint packet has an integrable
exponential envelope.
The critical-line Gram blocks converge to their nonsingular Hilbert
matrices, and distinct critical ordinates decouple by integration by
parts.  Thus the full finite Gram matrix has an invertible limit.
It follows that the orthogonal projection $P_{F,L}$ onto the finite
packet satisfies
\[
                    \|Q_\Lambda P_{F,L}\|_2=o(1).
\]

Since $0\leq D_L\leq S_L$, for every unit vector $u$ in the packet
\[
 \|(D_L^{1/2}-1)u\|^2
      \leq\langle u,(1-D_L)u\rangle=\|Q_\Lambda u\|^2.
\]
Applying the finite Bessel inequality to
$C_a^*D_L^{1/2}$ gives
\[
 \liminf\operatorname{Tr}(D_LC_aC_a^*)
 \geq\lim\operatorname{Tr}(P_{F,L}C_aC_a^*).
\]
The critical-line blocks contribute $m_\rho|a(\gamma)|^2$:
translation changes their slowly varying normalized polynomials by
$o(1)$, so $C_a^*$ acts asymptotically by the scalar
$\overline{a(\gamma)}$.
The endpoint blocks are finite Hardy model spaces.  An elementary
Malmquist orthonormal basis for a packet of upper-endpoint exponents
has Fourier transforms
\[
 \frac{\sqrt{2\alpha_j}}{\alpha_j+i(\tau-\gamma_j)}
 \prod_{k<j}\frac{\tau-\gamma_k+i\alpha_k}
                     {\tau-\gamma_k-i\alpha_k},
 \qquad\alpha_j>0,
\]
with each exponent repeated according to multiplicity.
The successive factors have modulus one on the real line.
Orthogonality follows by closing the rational inner-product integral
in the half-plane containing no pole after cancellation; completeness
follows from partial fractions and the dimension of the packet.
Reflection gives the lower-endpoint basis.
Their squared Fourier norms therefore sum to the Poisson kernels
in the statement, with multiplicity.  This proves the finite-packet
version of \eqref{eq:allzero-lower}.

To pass to all zeros, take increasing finite packets and use monotone
convergence.  This passage does not require a uniform Gram estimate
in the packet size.  For any Schwartz function $b$ and $0<a\leq1/2$,
split the Poisson integral at $|\tau-\gamma|=|\gamma|/2$ to obtain
\[
 \int P_a(\tau-\gamma)|b(\tau)|\,d\tau
                         \leq C_b(1+|\gamma|)^{-2}.
\]
The usual zero count $N_\chi(T)=O_\chi(T\log(2+T))$ makes the
right side summable over the zeros.  The critical-line part is
summable as well.  Compactly supported squared tests imply the
asserted domination for any vague subsequential limit.
Finally \eqref{eq:alljet-trace-input} transfers the inequality to
the uncompressed trace on the specified tests.
\end{proof}

\paragraph{What this proves.}
The result identifies the full unavoidable lower measure, including
the off-line zeros and every multiplicity.  It does not prove that
the defect measures are locally bounded, that they converge, or
that no extra positive measure remains.  Those are upper-bound
questions.
